\clearpage
\section{Concrete search for the copying UAM}
\label{sec:copying-search-phase}
This milestone defines the intrinsic search interactions on the copying graph
hierarchy and proves their operational laws. It covers the operation-search
rules of Figure~9 and the additional term-contraction and instance interactions
of Figure~18 (PDF p.~37). It uses the complete example operation alphabet.
Copying itself (Figure~17), substitution of stable arguments, and the compute
rules, including Figure~19 store updates, remain separate obligations.

The representation is a token mode and a wire of an actual typed hypernet.
The graph remains fixed during a search transition. No copying transition is
encoded as a linear compute step. A rewrite token suspends this phase; it is
not a successful value and does not mean the full copying machine is stuck.
An eventual concrete specimen instance must connect these coordinates with
focused graph states and combine them with copy and compute transitions.

\Needspace{12\baselineskip}
\subsection{Incidence rules and determinism}

\termdef{copying-search-phase}
The argument count is zero for constants and lambda, one for unary operations,
and two for application and binary operations. In particular, lambda's thunk
output is not an evaluated argument. Primitive labels are injective at every
depth: at depth zero project from the lifted label; at successor depth use
injectivity of the primitive summand. A position is a pair $(m,v)$ of mode and
vertex.

\lean{Copying/Search}{9}{12}{Copying.Operation.arguments}

\lean{Copying/Search}{14}{18}{Copying.primitiveLabel_injective}

\lean{Copying/Search}{21}{21}{Search.Position}

Entry at an operation with no evaluated arguments returns its finishing
mode at the same wire; otherwise it searches target zero. Entry at either
input of a term contraction produces rewrite mode at that input. Entry at an
instance produces backtrack mode at the same wire. Atom contractions, weakening,
stores, focus edges and holes have no entry instruction. On return from
argument $i$, the resume function searches argument $i+1$ if it exists among
the evaluated arguments; otherwise it uses the operation's finishing mode.
The optional result makes an absent target an absent transition.

\lean{Copying/Search}{25}{31}{Search.enter}

\lean{Copying/Search}{35}{37}{Search.resume}

The transition relation requires an actual incident edge. An entry witness
gives its exact primitive label, the current vertex in its source list, and
the entry equation. A return witness gives an operation edge with singleton
source, an index strictly below its argument count, the current vertex at
that target index, and the resume equation. Thus neither a label guess nor an
arbitrary adjacency suffices. Both entry and resume commute with every vertex
map: optional list lookup commutes with mapping, and all other branches retain
the same mode and mapped vertex. Injectivity is unnecessary for these two laws.

\lean{Copying/Search}{40}{47}{Search.Step}

\lean{Copying/Search}{50}{55}{Search.enter_map}

\lean{Copying/Search}{57}{60}{Search.resume_map}

For determinism, compare two entry witnesses using uniqueness of the edge
sourcing a given vertex; label injectivity then identifies their instructions.
For two return witnesses, target uniqueness identifies the edge, singleton
sources identify the root, and the duplicate-free target list identifies the
argument index. The two instruction equations give equal successors. Entry
and return cannot overlap because their source modes differ. The result-uniqueness
lemma allows explicitly equal source positions; determinism is its reflexive case.

There is no search step out of rewrite mode. Backtrack at an external input
also has no step, since an input cannot be an edge target. Finally, two finite
search runs from the same source to normal endpoints have equal endpoints and
equal lengths, by the generic deterministic trace theorem.

\lean{Copying/Search}{63}{91}{Search.result_unique}

\lean{Copying/Search}{93}{94}{Search.deterministic}

\lean{Copying/Search}{96}{97}{Search.no_rewrite}

\lean{Copying/Search}{99}{103}{Search.final_no_step}

\lean{Copying/Search}{106}{109}{Search.terminal_run_unique}

\Needspace{12\baselineskip}
\subsection{Exact transport across graph representations}

\termdef{copying-search-transport}
A search transport consists of bijections on vertices and edges, preservation
and reflection of each primitive label, exact preservation of the ordered
source and target lists, and preservation and reflection of input status.
It may relate different depth allowances. No box-body equation is required by
this interface because search does not enter boxes. It is a sufficient
interface for search only, not a replacement for recursive graph isomorphism.
Its position bijection preserves modes. Its inverse uses the inverse carrier
bijections, the reversed primitive equivalence, and mapped-list inverse laws.

\lean{Copying/SearchTransport}{10}{19}{Search.Transport}

\lean{Copying/SearchTransport}{21}{25}{Search.Transport.positions}

\lean{Copying/SearchTransport}{27}{37}{Search.Transport.symm}

Transport a step by mapping its incident edge and vertices. List membership,
indexed lookup and singleton sources are preserved, and the entry/resume map
lemmas supply the result equation. Applying the inverse transport proves
reflection. Induction on a finite trace gives correspondence at exactly the
same number of steps. Normality follows in both directions by transporting any
hypothetical outgoing step back to the original graph.

\lean{Copying/SearchTransport}{40}{52}{Search.Transport.step}

\lean{Copying/SearchTransport}{55}{61}{Search.Transport.step_iff}

\lean{Copying/SearchTransport}{63}{65}{Search.Transport.trace_iff}

\lean{Copying/SearchTransport}{67}{75}{Search.Transport.terminal_iff}

Recursive label equivalence preserves and reflects equality with each
primitive label: primitive cases are literal equality and box cases cannot
be primitives. Hence a typed boundary isomorphism supplies a concrete search
transport using its carrier bijections and incidence laws. This proves search
invariance under the existing recursive isomorphism discipline without
identifying box permutation with ordered isomorphism.

\lean{Copying/SearchTransport}{79}{90}{Search.primitive_equiv}

\lean{Copying/SearchTransport}{93}{100}{Search.ofBoundaryIso}

\Needspace{12\baselineskip}
\subsection{Value recognition and wire typing}

An entry or resume result in search mode belongs to the interacting edge's
target list, by the successful optional lookup. Consequently no step produces
search mode at an external input. For a backtrack entry result, case analysis
on the primitive and operation proves that its vertex is the unchanged root
and its label is a stable primitive. Conversely, a return witness has an
argument index and therefore an active operation; its resume instruction
cannot produce backtrack mode.

\lean{Copying/SearchValues}{8}{21}{Search.enter_search_mem}

\lean{Copying/SearchValues}{23}{32}{Search.resume_search_mem}

\lean{Copying/SearchValues}{35}{40}{Search.no_initial_target}

\lean{Copying/SearchValues}{43}{48}{Search.enter_backtrack}

\lean{Copying/SearchValues}{51}{52}{Search.argument_active}

\lean{Copying/SearchValues}{55}{63}{Search.resume_not_backtrack}

\termdef{copying-search-value-recognition}
It follows that every transition producing backtrack mode identifies an
actual passive-operation or instance edge. Its source is exactly the singleton
containing the returned vertex: the typed source length is one, and source
membership identifies that unique element. This is operational recognition of
a stable primitive, not an assumed machine law.
An interface equality cast preserves every shallow label predicate by equality
elimination. Applied to the primitive corolla and the already proved stability
criterion, this shows that the recognised primitive, isolated as its own
corolla with the required input interface, is a stable graph. It does not
claim that an arbitrary surrounding graph is stable.

\lean{Copying/SearchValues}{67}{80}{Search.backtrack_stable_edge}

\lean{Copying/SearchValues}{83}{86}{Search.labels_cast}

\lean{Copying/SearchValues}{89}{94}{Search.recognized_primitive_stable}

Typing is checked against the actual signatures. Successful entry has
only term sources; given a term root, its result wire is term. Every evaluated
argument index selects a term wire in the operation output signature. Resume
therefore also preserves term type. These finite signature facts are proved
by cases on the complete operation and primitive alphabets. For an actual
step, map its instruction equation along the graph's vertex-label function
and use the source/target typing equations. Both endpoint vertices then have
term type. In particular, the instance's atom output and lambda's thunk output
cannot become searched arguments.

\lean{Copying/SearchTyping}{8}{14}{Search.enter_input_term}

\lean{Copying/SearchTyping}{16}{21}{Search.enter_output_term}

\lean{Copying/SearchTyping}{24}{29}{Search.argument_type}

\lean{Copying/SearchTyping}{31}{38}{Search.resume_output_term}

\lean{Copying/SearchTyping}{41}{69}{Search.endpoints_term}

\Needspace{12\baselineskip}
\subsection{Padding and shallow context insertion}

One recursive raise supplies an actual search transport: outer vertex and
edge carriers are unchanged, primitive labels remain the same primitive,
ordered incidence is unchanged, and input status is unchanged. This argument
uses the definition of padding, not determinism. Identity transports and
composition are built from the corresponding bijections; the primitive,
incidence and input equations compose. Iterate this construction for padding
by $n$, transport across a depth-index equality by equality elimination, and
compose the two for padding to an arbitrary common allowance. Step reflection
and exact trace correspondence follow from the transport theorems.

\lean{Copying/SearchPadding}{9}{26}{Search.padding}

\lean{Copying/SearchPadding}{30}{36}{Search.Transport.refl}

\lean{Copying/SearchPadding}{38}{49}{Search.Transport.trans}

\lean{Copying/SearchPadding}{53}{55}{Search.pad}

\lean{Copying/SearchPadding}{57}{58}{Search.castDepth}

\lean{Copying/SearchPadding}{60}{61}{Search.padTo}

\lean{Copying/SearchPadding}{63}{65}{Search.padding_step_iff}

\lean{Copying/SearchPadding}{67}{71}{Search.padTo_trace_iff}

A successful entry edge cannot be a hole: label injectivity would identify
its primitive with a hole, whose entry instruction is absent. When filling a
shallow hole, a host interaction therefore retains its edge; return interactions
also retain their operation edge. Map its vertices along the host attachment
map. The plugging incidence equations and the entry/resume map laws prove the
same transition in the filled graph. For an interaction inside the filling,
use the inserted edge and inserted-vertex map instead. Neither map is assumed
to be injective. Induction transports whole host or inserted traces without
changing their lengths. These are preservation laws for shallow plugging;
reflection, nested specimen classification and full-machine locality remain
additional obligations.

\lean{Copying/SearchLocality}{8}{16}{Search.search_ne_hole}

\lean{Copying/SearchLocality}{20}{42}{Search.plug_host}

\lean{Copying/SearchLocality}{45}{61}{Search.plug_insert}

\lean{Copying/SearchLocality}{64}{71}{Search.plug_host_trace}

\lean{Copying/SearchLocality}{73}{79}{Search.plug_insert_trace}

\Needspace{12\baselineskip}
\subsection{The phase machine and exact observations}

\termdef{copying-search-completion}
Successful phase completion means backtrack mode at an external input. The
phase machine has exactly the search relation as its step relation and this
completion predicate as finality. Final normality is the input obstruction
proved above; determinism is the concrete search theorem. Generic counted
machine runs and graph traces coincide by induction on their constructors.

A transport preserves completion because it preserves mode and input status.
Combining this with exact trace transport gives exact machine-run correspondence.
For existential successful runs, an arbitrary target endpoint is pulled back
through the inverse position bijection; conversely map the original endpoint
forward. Thus successful runs exist at cost $n$ on one side exactly when they
do on the other. Instantiating the transport by arbitrary padding proves the
corresponding operational padding theorem for this phase. No full copying
execution instance is manufactured from it.

\lean{Copying/SearchMachine}{10}{11}{Search.Completed}

\lean{Copying/SearchMachine}{14}{21}{Search.machine}

\lean{Copying/SearchMachine}{23}{24}{Search.machine_deterministic}

\lean{Copying/SearchMachine}{27}{40}{Search.steps_iff_trace}

\lean{Copying/SearchMachine}{42}{43}{Search.Transport.completed_iff}

\lean{Copying/SearchMachine}{45}{47}{Search.Transport.machine_steps_iff}

\lean{Copying/SearchMachine}{50}{61}{Search.Transport.completes_at_iff}

\lean{Copying/SearchMachine}{65}{69}{Search.padding_completion_iff}

\Needspace{12\baselineskip}
\subsection{Regressions on the original translated graphs}

The term-contraction and instance rules specialise the entry constructor
with the exact Figure~18 instructions. A no-entry lemma rules out a competing
witness: source uniqueness identifies its edge and primitive-label injectivity
identifies its instruction, contradicting absence of an entry result. It
excludes search at atom contractions and stores. Concrete corolla regressions
cover both ordered term-contraction inputs and the instance input at every
depth, using their explicit source lists.

\lean{Copying/SearchExamples}{10}{12}{Search.term_contraction}

\lean{Copying/SearchExamples}{15}{17}{Search.instance_return}

\lean{Copying/SearchExamples}{20}{31}{Search.no_entry}

\lean{Copying/SearchExamples}{33}{35}{Search.atom_contraction_no_search}

\lean{Copying/SearchExamples}{37}{39}{Search.store_no_search}

\lean{Copying/SearchExamples}{42}{45}{Search.contraction_ports}

\lean{Copying/SearchExamples}{47}{50}{Search.instance_corolla}

For any translated abstraction body, the outer lambda edge is an actual
edge of the two serial assemblies. Its source is the external input; its
zero argument count returns backtrack there. The original box and single-use
distributors remain in place. This applies to the unchanged right-hand side
of Proposition~12.2 for every environment size $(k,h)$. Its one-step trace
ends at a normal input, so deterministic terminal-run uniqueness proves that
every successfully completed search run from this initial state has exactly
one step and that same endpoint. This says nothing yet about callbacks or
repeated use of the returned closure.

\lean{Copying/SearchExamples}{53}{58}{Search.abstraction_returns}

\lean{Copying/SearchExamples}{61}{63}{Search.constantResult_search}

\lean{Copying/SearchExamples}{66}{73}{Search.constantResult_exact_cost}

For an ordered multi-argument test, take the literal Figure~16 translation
of $1+2$ in the empty environment. The graph definition retains all
distributors. Its named head, left constant and right constant edges are the
corresponding finite assembly coordinates; argument $i$ is target $i$ of the
head, rather than a separately invented vertex.

\lean{Copying/SearchTraceExamples}{9}{10}{Search.additionGraph}

\lean{Copying/SearchTraceExamples}{12}{12}{Search.additionHead}

\lean{Copying/SearchTraceExamples}{13}{13}{Search.additionLeft}

\lean{Copying/SearchTraceExamples}{14}{14}{Search.additionRight}

\lean{Copying/SearchTraceExamples}{15}{16}{Search.additionArgument}

The five checked transitions are: enter the head and search the left
constant; return from that constant; resume at the head and search the right
constant; return from it; finish the head in rewrite mode at the root. Each
proof gives the actual edge, primitive label, ordered source and target
positions, and reduced instruction equation. The trace concatenates these
five witnesses. Since rewrite mode has no outgoing search step, every normal
search run from the same initial state has this endpoint and length five.
The padding trace theorem transports this concrete run to every larger depth
allowance. No arithmetic computation or store mutation is claimed by the test.

\lean{Copying/SearchTraceExamples}{19}{22}{Search.addition_enter}

\lean{Copying/SearchTraceExamples}{25}{28}{Search.addition_left}

\lean{Copying/SearchTraceExamples}{31}{32}{Search.addition_next}

\lean{Copying/SearchTraceExamples}{35}{38}{Search.addition_right}

\lean{Copying/SearchTraceExamples}{41}{42}{Search.addition_finish}

\lean{Copying/SearchTraceExamples}{45}{47}{Search.addition_trace}

\lean{Copying/SearchTraceExamples}{50}{53}{Search.addition_exact_cost}

\lean{Copying/SearchTraceExamples}{56}{60}{Search.addition_padded}

The full build checks 298 jobs. All three axiom audits pass: six abstract,
477 graph and 428 copying declarations, with only the standard logical axioms
\texttt{propext}, \texttt{Classical.choice} and \texttt{Quot.sound} where used.
Appendix~B removes the completed search-definition obligation while retaining
the copy, compute, refocusing and full specimen-interface tasks. Estimated
completion of Proposition~12.2 is about 55\% by mathematical effort, with a
broad 40--65\% uncertainty range (26 October 2026). This is a judgement about
the remaining mathematics, not a ratio of declarations or pages; the template
robustness proofs still carry substantial uncertainty.

